\documentclass{ltjsbook}
\usepackage{docmute} 
\usepackage{../../myStyle} 

\begin{document}
% \chapter{双対尺度法}
\chapter{Dual Scaling Method}
\label{x14_DualScaling}


% さて前回SEMを学んだことで，分散共分散行列をとことんまで分析し尽くす方法が手に入ったことになります。SEMは統合的な表現方法ですから，観測変数でも潜在変数でもいいですし，\keytermL{順序尺度水準}{じゅんじょしゃくどすいじゅん}以上の大小関係が表現できる数値データであれば，あらゆる表現ができるわけです。
Now, having learned about SEM in our last session, we now have a way to thoroughly analyze the variance-covariance matrix. SEM is an comprehensive way of representation, so it's okay whether it is an observed variable or a latent variable. If it's numerical data that can express a relationship of greater or lesser than the \keytermL{ordinal scale level}, then all kinds of representations are possible.

% ではこれで多変量解析はすべて理解したことになるのでしょうか。いえ，もちろんそうではありません。分散共分散行列の分解はできるようになりましたが，変数間関係の表現は分散共分散行列だけではありません。
Does this mean we've completely understood multivariate analysis? No, of course not. We can now decompose the variance-covariance matrix, but the representation of relationships between variables is not limited to just the variance-covariance matrix.
% ここまで，データの持つ情報は分散がすべて，変数間関係は共分散で表されるから分散共分散行列がすべてだ，と言わんばかりに話をしてきましたが，その枠を外すとどうなるでしょうか。
Up to now, we've been discussing as if all the information in the data is represented by the variance, and the relationships between variables are represented by the covariance, hence the covariance matrix is everything. But what happens if we move beyond that framework?

% \section{直線的ではない関係}
\section{Non-Linear Relationships}

% ご存知の通り，分散共分散行列を標準化した行列は相関行列といいますが，相関係数が$1.0$(あるいは$-1.0$)の状態を考えれば明らかなように，分散共分散行列(や相関行列)で表現されるのは変数の直線的な関係性に限った話です。心理学の場合は中庸が良いようなシーンも少なくありません。たとえば血圧と健康度の関係で言えば，低血圧でも高血圧でも不健康であり，ちょうどいい血圧が一番健康的だというのはすぐにわかります。このように中庸が良い場合は，散布図がU字型に現れてきます。散布図がU字型であれば，相関係数としては0近くになりますが，血圧と健康が無関係だとは誰も言えないでしょう。
As you know, the matrix obtained by standardizing the variance-covariance matrix is called the correlation matrix. However, as is clear when considering the state where the correlation coefficient is $1.0$ (or $-1.0$), the variance-covariance matrix (and the correlation matrix) can only express linear relationships between variables. In psychology, there are often instances where moderation is best. For instance, it's easy to understand that in the relationship between blood pressure and health, being either hypotensive or hypertensive is unhealthy, and having just the right blood pressure is the healthiest. In such cases where moderation is best, scatter diagrams appear in a U-shape. If the scatter diagram is U-shaped, the correlation coefficient will be close to 0, but no one can say that there is no relationship between blood pressure and health.


% ごくシンプルかつ具体的な例をあげてみましょう。血圧と頭痛の頻度について調査したとします。血圧は高い，普通，低いの3段階。頭痛の頻度は「ない」、「たまに」、「ときどき」，「いつも」の4段階です。表\ref{tbl::14_01}のような結果を見ると，この集計表に直線的な関係は確かになさそうですね。でも関係ないわけではない，と思いませんか。
Let's give a very simple and specific example. Suppose we conducted a survey on blood pressure and the frequency of headaches. Blood pressure is classified into three levels: high, normal, and low. The frequency of headaches is divided into four levels: "none," "occasionally," "sometimes," and "always". Looking at the results like in Table \ref{tbl::14_01}, there doesn't seem to be a linear relationship in this summary table. But don't you think it's not unrelated?

\begin{table}[ht]
	\begin{center}
		\caption{血圧と頭痛}
		\begin{tabular}{c|cccc}
			血圧 & ない & たまに & ときどき & いつも \\
			\hline
			高い & 0  & 0   & 3    & 2   \\
			普通 & 5  & 0   & 0    & 0   \\
			低い & 0  & 2   & 3    & 1   \\
		\end{tabular}
		\label{tbl::14_01}
	\end{center}
\end{table}


% このような場合は、「血圧が高い人か低い人は、頭痛がある」という傾向が見て取れるはずです。しかしこのデータ，機械的に血圧が高いを1，普通を2，低いを3とし，同様に頭痛もない〜いつもを1,2,3,4として相関係数を計算すると，$r=0.165$になります。相関関係では傾向をうまく読み取れていないのです\footnote{相関係数でお話ししましたが，分散共分散行列でも同じです。共分散で表現すると単位のせいで関係が直接的には理解できません。}。
In such cases, there should be a tendency for people with high or low blood pressure to have headaches. However, with this data, if we mechanically assign high blood pressure a 1, normal a 2, and low a 3, and similarly assign no headache to always a 1,2,3,4, and calculate the correlation coefficient, we get $r=0.165$. The correlation relationship is not correctly capturing the trend\footnote{I explained it using correlation coefficient, but the same is true for a covariance matrix. With covariance, we cannot directly understand the relationship due to units.}.

% その根本的な原因は，もちろんスコアリングの方法にあります。共分散を計算するには間隔尺度水準以上の情報が必要で，今回のような3,4件法では相関係数を計算して良い数字ではありません。
The fundamental reason, of course, lies in the method of scoring. To calculate the covariance, information above the interval scale level is required, and in cases like this time, it is not appropriate to calculate the correlation coefficient using a 3 or 4 item method.
% ではポリコリック相関係数のように順序尺度水準の相関係数なら良いか，といいたいところですが，それでもまだ不十分です。なぜなら，血圧が高い方から低い方まで，順番が保存されてしまっているからです。
I'd like to say that it would be okay if we had a correlation coefficient at the ordinal scale level, like the polychoric correlation coefficient. However, this is still insufficient. This is because the order is still preserved, from high blood pressure to low.

% 関係を導き出すには抜本的な改革が必要です。たとえば表\ref{tbl::14_02}のようにすればどうでしょうか。こうすると左下から右上にかけての直線的な関係がまだ見えてきます。
Radical reforms are necessary to derive relationships. For example, how about doing it like Table \ref{tbl::14_02}? In doing so, we can still see the linear relationship from the bottom left to the top right.
% これは血圧の順番を「高い」「低い」「普通」に並べ替えたものです。
This is the blood pressure sorted in the order of "high", "low", and "average".
% また，頭痛の順番も「いつも」と「ときどき」をひっくり返しました。
Also, the order of headaches was reversed from "always" to "sometimes".
\begin{table}[ht]
	\begin{center}
		\caption{並び替えられた「血圧と頭痛」}
		\begin{tabular}{c|cccc}
			血圧 & ない & たまに & いつも & ときどき \\
			\hline
			高い & 0  & 0   & 2   & 3    \\
			低い & 0  & 2   & 1   & 3    \\
			普通 & 5  & 0   & 0   & 0    \\
		\end{tabular}
		\label{tbl::14_02}
	\end{center}
\end{table}

% 順番を並び替えてしまったというのは，尺度水準的には数字を無視して「高い」「低い」「普通」というカテゴリーとして扱ったということになります。つまり\keytermL{名義尺度水準}{めいぎしゃくどすいじゅん}レベルにまで落として考えたのです。
The fact that the order was rearranged means that numerically, the categories were treated as "high", "low", and "normal" ignoring the numbers. In other words, it was considered at the \keytermL{Nominal scale level}.

% このように並べ替えて，血圧のスコアを高い$\to3$，低い$\to2$，普通$\to1$とし，また頭痛の頻度をいつも$\to3$，ときどき$\to4$と付け替えて相関係数を計算すると，今度は$r=0.810$になりました。こちらの方が線形性は高いことが数字でも確認できました。
After rearranging in this way, blood pressure scores are set as high $\to3$, low $\to2$, normal $\to1$, and the frequency of headaches is changed to always $\to3$, sometimes $\to4$. By calculating the correlation coefficient like this, it has now become $r=0.810$. This confirmed numerically that this method shows a higher degree of linearity.

% ここで行ったのはこのように，すべての反応カテゴリをただの言葉だと考えて，付与された数字を無視して並べ替えたことになります。そうすることで，左下から右上にかけて線形性を高めることができました。しかし「高い」と「低い」，「低い」と「普通」の配置も等間隔である必要はなく，もっと直線性がはっきりするように配置してやってもよいのでは，というアイデアが浮かびます。
Here, what we did was consider all response categories as just words, and rearrange them ignoring the numbers assigned. By doing so, we were able to increase the linearity from the lower left to the upper right. However, there is no need for "high" and "low", "low" and "normal" to be equally spaced; the idea comes up that it might be better to arrange them so that the linearity becomes more apparent.

% ここで考え方が反転したことに注意してください。一般的なリッカート尺度では，反応カテゴリに(シグマ法などで)数字をつけて，変数間の線形性を算出して意味を考えるのでした。
Please note that the way of thinking has been reversed here. In a typical Likert scale, we assign numbers to response categories (using methods like sigma), and we calculate the linearity between variables to interpret the meaning.
% ここでは逆に，変数間の線形性を最大にするように反応カテゴリに数字をつけてやろう，という考え方です。データの直線性というのはデータの特徴を最も強調し解釈しやすい形です。そのようにデータを整えるためには，反応カテゴリにどういう数字を付与すれば良いか，と考えるのです。この方法を\keyterm{数量化の理論}{すうりょうかのりろん}{Quantification Methods}といいます。
Here on the contrary, the idea is to assign numbers to the response categories in a way that maximizes the linearity between variables. Linearity of data highlights the features of the data in the most interpretable way. To organize the data in this way, you need to think about what numbers should be assigned to the response categories. This method is called Quantification Methods.

% 反応カテゴリに数字を与えるとき，名義尺度水準にまで落として，すなわち数字の意味を無くして並べ替えるような作業をします。もちろん分析対象が，初めから名義尺度水準であっても構いません。たとえば出身県を調査したようなデータがあったとして，他の変数との線形関係を最大にするように並べ替え，数字を付与することもできます。数量化の考え方は，名義尺度水準の変数に解釈しやすい数値を与えることとも言えるのです。
When giving numbers to response categories, we downscale to nominal scale level, or in other words, we work to eliminate the meaning of numbers and reorder them. Of course, it is okay if the subject being analyzed initially has nominal scale level. For example, if there is data obtained from a survey about the prefecture of origin, we can reorder to maximize their linear relation with other variables and assign numbers accordingly. The idea of quantification can also be said to provide easy-to-understand numbers to variables at nominal scale level.

% \section{林の数量化理論}
\section{Quantization Theory of Forests}
% 数量化の研究をしたのは，日本の偉大な統計学者，林知己夫(はやし ちきお)\footnote{林 知己夫(はやし ちきお、1918年6月7日 - 2002年8月6日)は、日本の統計学者。正四位勲二等、理学博士。統計数理研究所第7代所長。 社会調査・世論調査におけるサンプリング方法の確立を始め、数量化理論(Hayashi's Quantification Methods)の開発とその応用で知られる。1990年代以降、データの科学(Data Science)を提唱し、その研究・思想は現在へと引き継がれている。Wikipediaより。}という人です。その名を冠して林の数量化理論と呼ばれることがあります。
The person who did the research on quantification is a great Japanese statistician, Chikio Hayashi\footnote{Chikio Hayashi (June 7, 1918 - August 6, 2002) was a Japanese statistician. He held the ranks of Shoshi and Kuninoto, and had a doctorate in science. He was the seventh director of the Institute of Statistical Mathematics. He is known for establishing sampling methods in social surveys and public opinion surveys, developing Hayashi's Quantification Methods, and applying them. From the 1990s onward, he advocated data science, and his research and ideas have been handed down to the present day. From Wikipedia.}. It is sometimes referred to as Hayashi's Quantification Theory, named after him.

% 林はさまざまな研究成果を挙げており，論文ごとにデータに必要な分析方法を開発するというようなスタイルでした。その膨大な研究業績を，弟子である飽戸弘\footnote{飽戸 弘(あくと ひろし、1935年3月14日 - )は、日本の社会学者、東京大学名誉教授。専門は、社会心理学、コミュニケーション論。Wikipediaより。}が分類して，同じような分析方法ごとにI類，II類，III類，IV類，と呼んでいきました。
Hayashi produced a variety of research results, following a style where he developed the necessary analysis method for the data in each paper. His vast research achievements were categorized by his disciple, Hiroshi Akuto\footnote{Hiroshi Akuto (born March 14, 1935) is a Japanese sociologist and honorary professor at the University of Tokyo. He specializes in social psychology and communication theory. From Wikipedia.}, who classified them into Type I, Type II, Type III, and Type IV, according to similar analysis methods.

\begin{table}[ht]
	\caption{林の数量化}
	\begin{tabular}{c|c|c|c|c}
		手法    & 外的基準 & データ     & 目的           & 関連する手法 \\
		\hline
		数量化Ⅰ類 & 量的変数 & 質的変数    & 外的基準の予測      & 重回帰分析  \\
		数量化Ⅱ類 & 質的変数 & 質的変数    & 外的基準の判別      & 判別分析   \\
		数量化Ⅲ類 & なし   & 質的変数    & 変数間の関係の要約と記述 & 正準相関分析 \\
		数量化Ⅳ類 & なし   & 対象間の類似度 & 対象間の関係の要約と記述 & 多次元尺度法 \\
	\end{tabular}
	\label{tbl::14_03}
\end{table}


% 表\ref{tbl::14_03}にあるように，基本的には質的変数，すなわち名義尺度水準や\keytermL{順序尺度水準}{じゅんじょしゃくどすいじゅん}のデータが得られたときに，それを解釈するためにはどのような数値を割り振ってやれば良いか，という発想から生まれたものになっています。
As shown in Table \ref{tbl::14_03}, basically, it is something that was born from the idea of what kind of numbers should be assigned to interpret the data obtained at the nominal scale level and the ordinal scale level, i.e., qualitative variables.

% さきほどの表\ref{tbl::14_01}のようなデータの並べ替えについては，数量化でいうところのIII類に該当します。ただこの数量化III類はおもしろいもので，同時期に独立にこの手法がフランス，カナダでも発展しており，2つの別名を持っています。フランス学派が開発した手法は\keyterm{対応分析}{たいおうぶんせき}{correspondence analysis}といい，カナダ在住の日本人，西里静彦\footnote{西里 静彦(にしさと しずひこ、1935年6月9日 - )は、カナダの行動計量学者（計量心理学）。学位はPh.D.（ノースカロライナ大学・1966年）。トロント大学名誉教授、アメリカ統計学会フェロー、日本行動計量学会名誉会員。北海道十勝郡浦幌町出身(札幌市生まれ)。Wikipediaより。}が開発した手法は\keyterm{双対尺度法}{そうついしゃくどほう}{Dual Scaling}といいます。どの分析も，基本的にはカテゴリカルな変数について直線性を最大にする値を割り振ることを目的にします。
The rearrangement of data as in table \ref{tbl::14_01} corresponds to what is referred to as Type III in quantification theory. However, this Type III quantification is interesting in that it has developed independently in France and Canada during the same period and has two different names. The method developed by the French school is called correspondence analysis, and the method developed by Shizuhiko Nishisato, a Japanese person residing in Canada, is called dual scaling. Both analyses aim to assign values that maximize linearity for categorical variables. 

Note: Shizuhiko Nishisato (born June 9, 1935) is a Canadian behavioral metrician (quantitative psychologist). He has a Ph.D. from the University of North Carolina (1966). He is a Professor Emeritus at the University of Toronto, a Fellow of the American Statistical Association, and an Honorary Member of the Japanese Society of Behavioural Metrics. He was born in Sapporo and hails from the town of Urahoro in the Tokachi district of Hokkaido (Quoted from Wikipedia).

% カテゴリカル変数なので，分析の応用範囲は多岐にわたります。どのような変数でも名義尺度水準に落とすことができるからです。表\ref{tbl::14_04}には一般的なデータセットの形を示しました。IDがあって，Q1,Q2...と項目が列方向に並ぶ形です。一行が一人の反応を表しています。Q1がたとえば性別などの名義尺度水準の変数であっても，男性$\to1$，女性$\to2$のようにコード化するルールを決めて入力します。Q2はたとえばリッカート法で当てはまる，やや当てはまる・・・などの5件法だったとしましょう。その場合はシグマ法によるスコアを入れる，あるいは簡便的に当てはまるを5，やや当てはまるを4，といったように数字を割り振って入力しますね。
Since it is a categorical variable, the application range of analysis varies widely. This is because any variable can be reduced to the level of nominal scale. Table \ref{tbl::14_04} shows the general form of a dataset. It has an ID, and items such as Q1, Q2... are lined up in the column direction. Each row represents the response of one person. Even if Q1 is a variable at the level of nominal scale, such as sex, you decide on a rule to code it as men $\to1$. Women $\to2$ and input it. Suppose Q2 is a 5-point method such as "applicable", "somewhat applicable" etc. in the Likert method. In that case, you insert a score by the sigma method, or simply assign numbers such as 5 for applicable, 4 for somewhat applicable, etc, and input them.


% これをカテゴリカル変数と見なしてデータセットにした例が表\ref{tbl::14_05}です。IDは分析対象ではありませんから横に置くとして，Q1が2列に増えています。ID=1の人は男性なのですが，この人はQ1:Male=1かつQ1:Female=0というようにコード化されています。同様に，Q2には「どちらとも言えない」を選択しているのですが，これを3とするのではなく\texttt{00100}と5列にわたってコード化しているのです。このようにすることで，すべての変数をカテゴリーとして扱うことができるようになります。
This is an example of viewing this as a categorical variable and creating a dataset, as shown in Table \ref{tbl::14_05}. Since the ID is not the target of analysis, it is placed horizontally, and Q1 is increased to two columns. The person with ID=1 is a male, but this person is coded as Q1:Male=1 and Q1:Female=0. Similarly, for Q2, they have selected "neither can be said", but this is not considered as 3, but is coded as \texttt{00100} over 5 columns. By doing this, all variables can be treated as categories.


\begin{table}[ht]
    \begin{minipage}[t]{.30\textwidth}
        \begin{center}
            \caption{一般的なデータセット}
            \begin{tabular}{cccc}
                ID       & Q1       & Q2       & $\cdots$ \\\hline
                1        & 1        & 3        & $\cdots$ \\
                2        & 2        & 4        & $\cdots$ \\
                $\vdots$ & $\vdots$ & $\vdots$ &          \\\hline
            \end{tabular}
        \end{center}
        \label{tbl::14_04}
    \end{minipage}
    %
    \hfill
    %
    \begin{minipage}[t]{.60\textwidth}
        \begin{center}
            \caption{カテゴリカル化したデータセット}
            \begin{tabular}{ccccccc}
                ID       & Q1:M      & Q1:F     & Q2:5 & Q2:4 & Q2:3 & $\cdots$ \\\hline
                1        & 1         & 0        & 0    & 0    & 1    & $\cdots$ \\
                2        & 0         & 1        & 0    & 1    & 0    & $\cdots$ \\
                $\vdots$ & $\vdots $ & $\vdots$ &                               \\\hline
            \end{tabular}
        \end{center}
        \label{tbl::14_05}
    \end{minipage}
\end{table}


% また表\ref{tbl::14_01}を並べ替えた時のように，こうした名義尺度のデータの線形性が最大になるように，行だけでなく，列も並べ替えます。データの中で線形性が最大になるように，反応カテゴリ\textbf{と}回答者に数字を割り振るのです。
As when rearranging table \ref{tbl::14_01}, not only the rows, but also the columns are rearranged so that the linearity of such nominal scale data is maximized. Numbers are allocated to the reaction category \textbf{and} the respondents so that the linearity in the data is maximized.
% ちなみに表\ref{tbl::14_01}は集計されたデータであり，今回の表\ref{tbl::14_05}は集計前のデータです。カテゴリカルな変数の分析の場合は，どちらでも良いのです。行と列の関係が表されている数字であれば，「ある/ない」のような二値反応でも，集計された度数でも構いません。さらにいえば，行と列の関係が記述されていればなんでもいいのです。
By the way, Table \ref{tbl::14_01} is aggregated data, and Table \ref{tbl::14_05} is pre-aggregated data. In the case of analysis of categorical variables, either one is fine. Whether the numbers representing the relationship between rows and columns are binary responses like "Exist/Not exist", or aggregated frequencies, it doesn't matter. Further speaking, as long as the relationship between rows and columns is described, anything is acceptable.

% これまでの回帰分析，因子分析からSEMに至るまでの流れは，変数間関係だけを考えてきました。N行M列のデータ行列の計算の途中で，行に関する情報は平均化して潰されてしまい，最終的には$M \times M$サイズの正方行列だけを扱うことになったのでした。そしてM個の変数に\keytermL{因子負荷量}{いんしふかりょう}などの重みをつけて考察してきました。
The flow from regression analysis and factor analysis to SEM has only considered the relationship between variables. In the process of calculating the N-row M-column data matrix, information about the rows is averaged out and eventually only a square matrix of size M x M is handled. Then, weights such as factor loadings were applied to the M variables and examined.
% 今回のカテゴリカルなデータは，$N\times M$行の矩形行列をそのまま分析し，行と列の両方に重みをつけます。正方行列でも矩形行列でも，変数と変数，変数と回答者がどのように関係しているかというところを見るという意味では同じです。SEMでは分散共分散行列や相関行列が，双対尺度法ではクロス集計表や素データがその分析対象になります。数学的な面から説明すると，正方行列の場合は\keytermL{固有値分解}{こゆうちぶんかい}をして，\keytermL{固有値}{こゆうち}と\keytermL{固有ベクトル}{こゆうべくとる}を求め，固有ベクトルが変数の重みになるのでした。
The categorical data this time analyzes the $N\times M$ rectangular matrix as it is, and assigns weights to both rows and columns. Whether it's a square matrix or a rectangular matrix, it's the same in terms of looking at how variables and respondents are related. In SEM, the variance-covariance matrix or the correlation matrix becomes the analysis target, while in the dual scaling method, the cross tabulation table or raw data is the analysis target. If explained from a mathematical aspect, in case of a square matrix, we perform eigenvalue decomposition to obtain eigenvalues and eigenvectors, and the eigenvectors become the weights of the variables.
% 矩形行列の場合は\keyterm{特異値分解}{とくいちぶんかい}{Singular value decomposition}と呼ばれ，行および列に対応する\keytermL{特異ベクトル}{とくいべくとる}をその重み，座標と考えることになります。本質的には同じようなものだと思っていただければと思います。
In the case of rectangular matrices, this is called \keyterm{Singular Value Decomposition}, and it is considered that the \keytermL{singular vectors} corresponding to rows and columns are its weights and coordinates. You can essentially think of them as the same thing.

% \section{双対尺度法による分析}
\section{Analysis by Dual Scaling Method}
% それでは実際の分析例を見てみましょう。次のコード\texttt{code:\ref{code:14_01}}は，\texttt{MASS}パッケージに含まれるサンプルデータ\texttt{caith}を対応分析で分析するものです。
Let's take a look at an actual analysis example. The following code \texttt{code:\ref{code:14_01}} analyzes the sample data \texttt{caith} included in the \texttt{MASS} package using correspondence analysis.

\begin{lstlisting}[language=c++,label={code:14_01},caption={双対尺度法による分析}]
library(MASS)
caith
result <- corresp(caith,nf=min(nrow(caith),ncol(caith)-1))
result
plot(result)
\end{lstlisting}


% \paragraph{コード解説}
\paragraph{Code Explanation}
\begin{description}
	\item[1行目] パッケージ\texttt{MASS}のよみこみ
	\item[2行目] サンプルデータ\texttt{caith}を表示させる
	\item[3行目] 対応分析関数\texttt{corresp}で分析した結果を\texttt{result}オブジェクトに入れる
	\item[4行目] 結果の出力
	\item[5行目] 結果のプロット
\end{description}


% データ\texttt{caith}はスコットランドのCaithness地方に住んでいた人に対してなされた調査で，髪の毛の色と目の色の関係を集計したものです。列方向に髪の毛の色，行方向に目の色が入っていますね(表\ref{tbl::14_06})。
The data \texttt{caith} is from a survey conducted among people living in the Caithness region of Scotland. It compiles the relationship between hair color and eye color. The hair color is arranged in columns and the eye color in rows (Table \ref{tbl::14_06}).

\begin{table}[ht]
	\caption{データ\texttt{caith}の中身}
	\centering
	\begin{tabular}{lrrrrr}
		\hline
		       & fair & red & medium & dark & black \\\hline
		\hline
		blue   & 326  & 38  & 241    & 110  & 3     \\
		light  & 688  & 116 & 584    & 188  & 4     \\
		medium & 343  & 84  & 909    & 412  & 26    \\
		dark   & 98   & 48  & 403    & 681  & 85    \\
		\hline
	\end{tabular}
	\label{tbl::14_06}
\end{table}


% この行・列のカテゴリにはとくに順序性などありませんが，分析することで最も線形性の高いウェイトをつけることができます。もちろん1次元で表現できない可能性があり，その場合は第二，第三と次元数を増やして行きます。データの行数あるいは列数の小さい方マイナス1次元まで求めることができます。それを表現しているのが，関数\texttt{corresp}のオプション\texttt{nf}で，行数\texttt{nrow(caith)}，列数\texttt{ncol(caith)}の小さい方(\texttt{min}関数)マイナス1を指定しています。
There is no particular order to the categories of these rows and columns, but by analyzing, we can apply the highest linearity weight. Of course, there may be a possibility that it cannot be represented in one dimension, in which case, we continue to increase the number of dimensions to the second and third dimensions. We can determine up to one dimension less than the fewer of the number of rows or columns of the data. This is represented by the option \texttt{nf} of the function \texttt{corresp}, where \texttt{nrow(caith)}, \texttt{ncol(caith)}, the smaller of the two (\texttt{min} function), minus one is specified.

% 結果は出力\ref{RS:out14_01}のようになります。行および列にスコアがついていますね。この値を尺度値として使うと，データの相関係数が最大になるというような指標化がなされたのです。
The result will look like Output \ref{RS:out14_01}. There are scores on rows and columns. These values are used as scale values, indicating that the data has been made into an index such that the correlation coefficient becomes maximum.
\begin{Rscreen}{対応分析の結果}{out14_01}
	\begin{verbatim}
First canonical correlation(s): 4.463684e-01 1.734554e-01 2.931691e-02 1.134031e-16 

 Row scores:
              [,1]       [,2]       [,3] [,4]
blue   -0.89679252  0.9536227  2.1884132    1
light  -0.98731818  0.5100045 -1.0837859    1
medium  0.07530627 -1.4124778  0.1894089    1
dark    1.57434710  0.7720361 -0.1482208    1

 Column scores:
              [,1]       [,2]       [,3]       [,4]
fair   -1.21871379  1.0022432  0.4271282 -0.8692696
red    -0.52257500  0.2783364 -4.0268545 -1.3400421
medium -0.09414671 -1.2009094  0.1103959 -0.8453208
dark    1.31888486  0.5992920  0.3450676 -1.2251588
black   2.45176017  1.6513565 -1.5736976  1.1609621
\end{verbatim}
\end{Rscreen}


% 結果はプロットされたものを見た方がわかりやすいかもしれません。図\ref{fig::14_01}にプロットを示しました。
The results might be easier to understand if you look at the plotted ones. The plot is shown in Figure \ref{fig::14_01}.
% たとえば横軸，dim1にそって目の色を見ていくと，light -- blue -- medium -- darkの順に並べた方が良い，ということがわかります。とくにlightとblueは近いのであまり大きな意味的違いはないことがわかります。
For example, if you look at the color of the eyes along the horizontal axis, dim1, you would see that it is best to arrange them in the order of light -- blue -- medium -- dark. In particular, light and blue are close, so you can see that there is not much significant semantic difference.
% 同様に髪の毛の色は，fair -- red --medium --dark --blackの順に並べられることになります。具体的な数字はそれぞれ出力の一列目の通りです。
Similarly, the colors of hair are arranged in the order of fair -- red -- medium -- dark -- black. The specific numbers are as listed in the first column of the output.
% このdim1だけでは表現できない違いが，dim2で表されており，それが図では上下の広がりとして示されていますね。
The difference that cannot be expressed with just dim1 is represented by dim2, and that is shown in the figures as the vertical spread.

% この並びが，新しく構成された次元だと考えることができます。
This order can be considered as a newly constituted dimension.
% そしてこれを見ると，たとえば髪の毛と目の色がどちらもdarkである場合，この二点は近くにプロットされています。この二点が近くにあるということは，両者の結びつきが強いということです。髪の毛が暗い人は目の色も暗い人が多い，といった関係の強さが見て取れます。
And when you look at this, for example, if both hair color and eye color are dark, these two points are plotted nearby. This proximity of the two points means that there is a strong connection between the two. You can see the strength of relationships such as that people with dark hair often have dark eyes as well.

\begin{figure}[ht]
    \begin{center}
        \includegraphics[width=0.8\linewidth,keepaspectratio]{../images/14_DualScaling/Rplot14_01.png}
        \caption{分析結果のプロット}
        \label{fig::14_01}
    \end{center}
\end{figure}


% ちなみにこれは対応分析独特の表示法で，行と列の変数が1つの画面にうつされていますね。
By the way, this is a unique representation of correspondence analysis, and both row and column variables are displayed on one screen.
% 厳密には，行ベクトルが作る次元，列ベクトルが作る次元は別物です。ですから，列の要素を行の空間に写像するためには変換が必要で，逆もまた然り，なのです。
Strictly speaking, the dimensions created by the row vectors and the column vectors are different. Therefore, transformations are needed to map the elements of a column into the space of a row, and vice versa.
% 双対尺度法の場合はこの違いを厳格に捉えます。対応分析の場合は，相互に写像しあった座標を一枚の図にすることで情報を圧縮しています。
In the case of dual scaling, this difference is strictly captured. In the case of correspondence analysis, information is compressed by mapping coordinates to each other and creating a single diagram.
% 数学的には同じモデルであっても，表現の仕方や表示の仕方に，モデル作成者の考え方が反映されているとも言えるでしょう。
Even for the same mathematical model, it could be said that the way of representation and display reflects the thoughts of the model creator.

% \section{テキストマイニングへの応用}
\section{Application to Text Mining}
% さて，本講で紹介した分析方法は，名義尺度水準を対象にしていて，かつ外的な基準がなく内的な構造を明らかにしようとするものだということがわかってきました。いわば名義尺度水準における因子分析ですね。
You should now understand from this lecture that the analytical method introduced is targeted at the nominal scale level and attempts to clarify the internal structure without an external benchmark. In other words, it's a factor analysis at the nominal scale level.

% 名義尺度水準を対象にしたということは，およそ言語化できたものはすべて分析の対象にできる，とも言えます。最も低次元で一対一対応した言葉の世界の数字だからです。
The fact that we targeted the nominal scale level means that we can say that everything that can be roughly verbalized can be a target for analysis. This is because it is a number in the world of words that corresponds one-to-one at the lowest dimension.


% 応用例として，\keyterm{テキストマイニング}{てきすとまいにんぐ}{Text Mining}を紹介しておきましょう。別名自然言語処理とも言われますが，これはテキスト，すなわち普段の「言葉」に潜む関係を分析する手法です\footnote{マイニングとは鉱脈を掘る，という意味です。}。
Let's introduce \keyterm{Text Mining}{てきすとまいにんぐ}{Text Mining} as an applied example. Also known as Natural Language Processing, this is a method of analyzing the relationships hidden in text, or the 'words' we use daily\footnote{Mining means to dig a vein.}.
% 文章を対象にしますから，小説や新聞記事などはもちろん，日記や逐語録なども分析の対象にしてしまおうというものです。心理学の分野では面接法などにおいて，クライアントがどのように語るかをすべて記録することがありますが，これをみて「うーん，こういうことを考えているんだな」と読み取るのは主観的な判断によるものです。ここにテキストマイニングを用いれば，機械的にどう言った言葉の使われ方をしているかを分析できます。あるいはツイッターやフェイスブックなどのSNSでどのような言葉，記事が流行しているかと言ったことを分析する，というような使い方もできます。分析対象が一気に広がりますね。
The text is targeted, so not only novels and newspaper articles, but also diaries and verbatim records are targeted for analysis. In the field of psychology, there are instances where everything a client says in an interview is documented, and interpreting this with thoughts like "Hmm, they're thinking this way" involves subjective judgment. If text mining is used here, you can mechanically analyze how words are being used. It's also possible to analyze what kind of words and articles are trending on social media platforms like Twitter and Facebook. The range of analysis definitely expands.

% このテキストマイニングがやっていることは二段回に分かれ，第一段階が\keyterm{形態素解析}{けいたいそかいせき}{morphological analysis}，第二段回が多変量解析です。第一段階は，「今日はいい天気ですね」といった平文を「きょう」「は」「いい」「てんき」「です」「ね」といった要素に分解することを指します。英語のような分かち書きがされる言語であればこれは簡単なのですが，日本語の場合は分かち書きされていないことに加え，漢字，ひらがな，カタカナなど表記方法もさまざまですから，この分析をするための特別な解析エンジンが必要です。幸い，すでに開発されているものがフリーソフトウェアとして利用できます。時間がかかりますがこれらを使うと，品詞ごとに分解し，その原形(活用する前の形)や活用形は何か，と言ったことを一覧にしてくれます。
This text mining is divided into two stages: the first stage is called morphological analysis, and the second stage is multivariate analysis. The first stage involves breaking down plain sentences such as "Today is nice weather" into elements like "today", "is", "nice", "weather", etc. This is easy in languages like English, which use separate words, but it is not so simple in Japanese where words are not separated and there are various forms of notation, such as Kanji, Hiragana, and Katakana. Therefore, a special analysis engine is required to perform this analysis. Fortunately, there are developed programs available as free software. Although it takes time, these can break down texts into parts of speech and list the original forms (before inflection) and inflected forms.

% 自然言語ですので大量の分割がなされますが，それを言葉同士の関係を表す行列の形で表現します。使われている品詞の原形ごとに誰が何回発言したかとか，各要素が1回の文章の中で同日に使われた回数をカウントするなどして，言葉と人，言葉と言葉の関係をデータ化するのです。
This is natural language, so a lot of segmentation is done, but it is expressed in the form of a matrix that indicates the relationship between words. By counting things like how many times each original form of the used part of speech was spoken, or how many times each element was used on the same day in a single sentence, we digitize the relationships between words, and between words and people.
% データ化できればあとはこっちのものです。数量化できるのですから，その言葉群のなかでどの言葉とどの言葉が近いのか，他のどの変数と関係するのか，と言った分析をできます。
Once it is digitized, the rest is ours. Since we can quantify it, we can analyze which words are close to each other in that group of words, and what other variables it is related to.

% テキストマイニングについてはRでもできますし，専門的なソフトウェアがあります。詳しくは\citet{Higuchu20200406}を参照してください。
You can do text mining in R, and there is also specialized software for it. For more details, please refer to \citet{Higuchu20200406}.


% \section{課題}
\section{Assignment}
% 今日の授業でおこなったすべての次の計算をするRコードを提出してください。ファイル形式はRスクリプトかRmdとします。なお提出されたコード単体でバグがなく動くことが確認できないものは，未提出扱いになります。コードの書き方などわからないところがあれば，曜日別TAか小杉までメールで連絡し，指導を受けてください。
Please submit the R code for all the calculations made in today's class. The file format should be an R script or Rmd. Please note that if the submitted code cannot be confirmed to run without bugs, it will be treated as not submitted. If you have any questions about how to write the code, please contact the weekday TA or Kosugi by email for guidance.


\end{document}
